\section{R15 执行记录：仅消除 $g_{\mathrm{total}}$ workspace 的融合试验}

\subsection{本轮目标与假设}

R14 的完整 workspace fusion 在 exactness 阶段失败，不能继续作为性能候选。
R15 将改动缩小到一个可以独立审计的数组：K1 原本把每个 tile 的
\(g_{\mathrm{total}}[D]\)（FP32）写入 workspace，K2 再通过 TMA 读取它。
候选分支不再写入和读取这一个数组，而是在 K2 读取原始 BF16 gate tile 后，
按原公式重新计算 gate、累计和以及 \(\exp_2\) 因子。其余五个 workspace、
TMA descriptor、线程组织和 recurrence 均保持不变。

该设计的预期收益是减少一组 global store/load；代价是 K2 增加 gate 的
转换、\(\tanh\) sigmoid、求和和指数计算。由于这两项变化位于不同 kernel，
不能从 DRAM 字节数直接推断端到端收益，必须用同一源码、同一设备和同一
CUDA-event harness 配对测量。

\subsection{实现与构建}

新增宏 \code{C1\_K1\_K2\_FUSED\_GTOTAL}，默认值为 0。K1 中该宏关闭时
仍写入原始 \(g_{\mathrm{total}}\)；打开时跳过该 TMA store。K2 中打开时
增加一个 BF16 raw-gate staging tile，由计算线程按列重建 \(g_{\mathrm{total}}\)。
没有修改 Python API，也没有改变 workspace ABI 的总分配接口；因此实验
分支可以和 baseline binary 交替安装。

远端构建使用 B300、\code{sm\_103a}、CUDA 13.0 和现有虚拟环境。候选与
control 分别从干净 checkout 构建，随后运行 smoke、extended exactness 和
配对 benchmark。轻量产物位于 \code{results/r15\_remote/}：
\code{r15\_gtotal\_smoke.json}、\code{r15\_gtotal\_extended.json}、
\code{r15\_gtotal\_candidate\_bench\_v2.json} 以及
\code{r15\_base\_bench.json}。

\subsection{正确性结果}

基础 smoke 的五个组合覆盖 \(T=16,17,64\)、\(H=1,96\)、BF16/FP32 state、
state in/out 和 no-state；extended 集合再覆盖 varlen、tail、\(B=2\) 和
FP32 output-only。候选的 output 与 final state 在全部十个组合中均为 exact：
\code{output\_different=0}、\code{state\_different=0}。这说明在当前
CHUNK=16、既定舍入顺序下，重建的终端 gate factor 与 K1 写回值一致，且
该局部融合没有改变 recurrence 状态。

\subsection{配对 benchmark}

下表给出 CUDA-event median；candidate/base 来自同一 B300 harness，比例定义
为 \(t_{\rm base}/t_{\rm cand}-1\)，正值才表示 candidate 更快。

\begin{center}
\small
\begin{tabular}{@{}l r r r@{}}
\toprule
case & candidate (ms) & base (ms) & candidate speedup \\
\midrule
fixed $T=16,H=1$   & 0.026048 & 0.026016 & $-0.12\%$ \\
fixed $T=16,H=4$   & 0.030080 & 0.030112 & $+0.11\%$ \\
fixed $T=16,H=96$  & 0.030880 & 0.030144 & $-2.38\%$ \\
fixed $T=8192,H=1$  & 1.878432 & 1.262880 & $-32.77\%$ \\
fixed $T=8192,H=4$  & 1.894848 & 1.284544 & $-32.21\%$ \\
fixed $T=8192,H=96$ & 2.339168 & 1.751584 & $-25.12\%$ \\
varlen $17+33,H=4$ & 0.042208 & 0.038144 & $-9.63\%$ \\
$B=2,T=64,H=4$   & 0.040320 & 0.036192 & $-10.24\%$ \\
\bottomrule
\end{tabular}
\end{center}

短序列的差异在测量噪声附近；长序列则稳定退化。候选节省了一个
\(H\,T_{\mathrm{tiles}}\,D\) 的 FP32 workspace 通路，却在每个 K2 tile
重复执行 gate 近似、累加和 \(\exp_2\)。在 \(T=8192,H=96\) 上，增加的
计算临界路径远大于一次 workspace TMA store/load 的成本，故总时间反而
增加约三分之一。这是“减少 DRAM traffic 必然加速”的反例。

\subsection{接入决策}

\begin{decision}{R15 接入决策}
R15 的局部融合通过了十个 exactness case，但没有一个长序列完整 output
case 超过同源 baseline；\(T=8192,H=96\) 候选慢约 33.55\%。因此
\code{C1\_K1\_K2\_FUSED\_GTOTAL} 保持 opt-in，默认值不变。若继续减少
workspace，应优先寻找能复用 K1 已计算的 gate 或真正重叠 producer/consumer
的 pipeline，而不是在 K2 重新计算相同非线性函数。
\end{decision}

\begin{record}{R15 结论}
R15 给出了一个干净的因果对照：只去掉一个 FP32 workspace 数组即可保持
完整 output/state exact，但由于重复 gate 计算，长序列端到端时间显著增加。
该结果支持“workspace 字节数减少并不等于 K2 加速”的判断，也为后续
workspace fusion 设定了约束：必须同时消除冗余计算，或将计算放入已有
producer pipeline 的空闲周期。
\end{record}
